\chapter{Guiamento híbrido PN/PP}
\label{cap:guiamento}

\section{Lei de controle}
Arquivo \code{control/hybrid\_guidance\_controller.cpp}. A aceleração é decomposta em
componente \textbf{lateral} (perpendicular à LOS, gira a LOS) e \textbf{radial} (ao
longo da LOS, controla a distância):
\begin{align}
  \vect a_{PN} &= N\,V_c\,(\vect\Omega\times\rhat),\quad
     \vect\Omega = \frac{\rr\times\dot\rr}{\lVert\rr\rVert^2},\quad V_c = -\dot r
     && \text{(navegação proporcional)}\\
  \vect a_{PP} &= K\,\max(\lVert\vect v_c\rVert,v_{min})\,
     \left(\rhat - (\rhat\cdot\hat{\vect v}_c)\hat{\vect v}_c\right)
     && \text{(perseguição pura)}\\
  \vect a_{lat} &= \left[\alpha\vect a_{PN}+(1-\alpha)\vect a_{PP}
     + k_{lat}\,\dot\rr_\perp\right]_\perp\\
  \vect a_r &= \left(k_r(r-d) + k_{rd}\,\dot r\right)\rhat .
\end{align}
\code{yaml}: $N=4$, $K=1$, $\alpha=0{,}5$, $k_{lat}=1$, $k_r=0{,}8$, $k_{rd}=1{,}2$,
$v_{min}=\SI{1}{m/s}$.

\section{Malha radial (distância)}
Ao longo da LOS, $\ddot r = a_{t,r} - a_{c,r}$. Com $a_{c,r}=k_r(r-d)+k_{rd}\dot r$:
\begin{equation}
  \ddot{\tilde r} + k_{rd}\dot{\tilde r} + k_r\tilde r = a_{t,r},\qquad \tilde r=r-d .
\end{equation}
É um PD clássico, então valem \eqref{eq:pd-polos} e o LGR. Com o \code{yaml}:
$\omega_n=\sqrt{0{,}8}=\SI{0,89}{rad/s}$ e $\zeta=1{,}2/(2\cdot0{,}89)=0{,}67$ ($t_s\approx\SI{6,7}{s}$,
mais lento que LQR/MPC). Diferentemente dos controladores de aceleração anteriores,
aqui $\dot r$ vem do KF (velocidade relativa \emph{medida}). O atraso entra também no
termo derivativo, então use $L(s)=(k_r+k_{rd}s)\,T_v(s)\,e^{-s\tau}/s^2$ na análise de margem.

\section{Malha lateral (linha de visada)}
Linearizando para pequenos ângulos de LOS $\lambda$ e alcance $r\approx d$:
$r\ddot\lambda \approx a_{t,\perp} - a_{c,\perp}$. Para a PN pura,
$a_{c,\perp}=N V_c\dot\lambda$. A teoria clássica de interceptação
(Zarchan) recomenda $N\in[3,5]$: com $N=3$ a demanda de aceleração é mínima contra um
alvo não manobrante, e $N>3$ antecipa manobras.

\begin{atencao}
A PN foi criada para \emph{interceptação} ($V_c>0$ grande). Na perseguição com standoff,
em regime $V_c\to0$ e o termo PN se anula. Quem mantém a LOS é o termo $k_{lat}\dot\rr_\perp$
(casamento de velocidade lateral) e o PP. Na fase de regime, portanto, a malha lateral
é aproximadamente $a_\perp=k_{lat}\dot\rr_\perp$: um amortecedor de 1ª ordem com
constante $1/k_{lat}$, sem realimentação de posição lateral. O ``erro lateral'' é
absorvido girando a LOS, e não há ponto lateral fixo a perseguir.
\end{atencao}

\section{Cálculo dos ganhos e ajuste ótimo}
\begin{enumerate}
  \item Radial: escolha $\omega_n,\zeta$ e faça $k_r=\omega_n^2$, $k_{rd}=2\zeta\omega_n$.
        Confira a margem com PX4 + atraso.
  \item Lateral: $N=3$--$4$; $k_{lat}\approx1/\tau_{lat}$ com $\tau_{lat}\sim$ 0,5--\SI{1}{s}.
  \item $\alpha$: compromisso entre PN (eficiente em aproximação) e PP (robusto a ruído
        de LOS, que a PN amplifica por depender de $\dot\lambda$).
  \item Ajuste fino: \code{otimizar\_ganhos.m} com $\theta=[N,\alpha,k_r,k_{rd},k_{lat}]$.
\end{enumerate}
